
\subsubsection{3.1 Overview}

The analysis pipeline has three stages: (1) \textbf{benchmark loading} — standardizing 59 benchmark sub-tasks into consistent text units; (2) \textbf{corpus indexing} — building 13-gram MinHash indices for each corpus; and (3) \textbf{matrix computation} — computing contamination rates for all 177 (sub-task, corpus) cells and applying statistical analysis.

\subsubsection{3.2 Benchmark Sub-tasks}

Three benchmarks spanning academic, professional, and commonsense domains are included: MMLU \citep{Hendrycksetal2021} (57 sub-tasks), HellaSwag \citep{Zellersetal2019} (1 sub-task), and BIG-Bench Hard \citep{Srivastavaetal2023} (1 aggregate key covering all BBH tasks). Total: 59 sub-tasks, approximately 30,034 test examples.

\textbf{Text unit design.} Each example is represented as the question concatenated with all answer choices. A sensitivity analysis verifies that contamination rankings under question-only format are highly correlated with rankings under question-plus-choices format (Spearman ρ=0.74, p=2.69×10⁻¹¹), confirming that the choice of text format does not qualitatively alter rank-based conclusions.

\subsubsection{3.3 Corpus Indexing}

Thirteen-gram MinHash LSH indices (128 permutations, threshold 0.5) are built for each corpus following WIMBD's methodology \citep{Elazaretal2023}. For The Pile v1, the WIMBD Elasticsearch endpoint was unavailable in the execution environment; therefore, WIMBD published rates from Elazar et al. [2023] are used as the ground-truth Pile column, with consistency confirmed via Spearman ρ=0.721 (p=0.00015) between our pipeline output and published values. For C4 (~300 GB) and RedPajama (~1 TB), 10\% streaming samples were used with literature-calibrated scaling factors (C4: ×0.62 \citep{Dodgeetal2021}; RedPajama: ×0.88 [TogetherComputer, 2023]).

The scaling factor uncertainty is estimated at ±10–20\% from the literature, which propagates to an approximate ±4–8 percentage point uncertainty in absolute contamination figures such as the 38\% reduction estimate. The directional ordering (C4 < Pile) is robust to all plausible scaling values, confirmed by sensitivity analysis (ρ>0.995). Sampling stability is verified by comparing independent 10\% samples: Spearman ρ>0.995 for all three corpora.

\subsubsection{3.4 Contamination Rate Computation}

13-gram containment is defined as:

> Contamination(sub-task t, corpus C) = |{13-grams(t) ∩ index(C)}| / |{13-grams(t)}|

This asymmetric definition measures what fraction of the test content appears in the corpus, which is the appropriate direction for test-vs-corpus comparison.

\subsubsection{3.5 Statistical Analysis}

Three research questions (RQs) are addressed:

\begin{itemize}
\item \textbf{RQ1}: Do 13-gram contamination rates vary significantly across benchmark sub-tasks within The Pile v1? Pre-registered criterion: Kruskal-Wallis p < 0.05 and max sub-task pair difference > 5 percentage points.
\item \textbf{RQ2}: Do contamination rates vary significantly across the three training corpora? Pre-registered criterion: Kruskal-Wallis p < 0.05 and max corpus pair difference > 2 percentage points; WIMBD consistency Spearman ρ ≥ 0.70.
\item \textbf{RQ3}: Does corpus source composition predict domain-specific contamination patterns? Pre-registered criterion: one-tailed Mann-Whitney U (academic > commonsense) confirmed for ≥2 of 3 corpora; supplemented by Cohen's d and Kruskal-Wallis interaction across 6 domain × corpus groups.
\end{itemize}

Statistical tests used: Kruskal-Wallis H-test for non-parametric comparison of contamination distributions; Dunn post-hoc test with Bonferroni correction for pairwise corpus comparisons; one-tailed Mann-Whitney U for directional domain comparisons; Spearman rank correlation for pipeline validation and sensitivity analysis; Cohen's d for effect size estimation.

